\documentclass[10pt,onecolumn]{article}
%% headers.tex — minimal preamble so paper.tex compiles standalone.
%%
%% When the paper is submitted through the JOSS workflow the Open Journals
%% bot supplies its own header/template and this file is ignored; it exists
%% so that `pdflatex paper.tex` works locally for length checking.

\usepackage[T1]{fontenc}
\usepackage[utf8]{inputenc}
\usepackage{lmodern}
% Asymmetric margins: the JOSS sidebar lives in a wide right margin,
% so the text block is narrowed rather than the marginpar clipped.
\usepackage[left=1.0in,right=2.05in,top=1.0in,bottom=1.0in,
            marginparwidth=1.5in,marginparsep=0.14in]{geometry}
\usepackage{amsmath,amssymb}
\usepackage{graphicx}
\usepackage{authblk}
\usepackage{xcolor}
\usepackage{listings}
\usepackage{fancyvrb}
\usepackage{caption}
\usepackage{enumitem}
\usepackage[colorlinks=true,allcolors=blue]{hyperref}

\definecolor{linky}{RGB}{0,110,190}
\definecolor{codebg}{RGB}{247,247,247}
\definecolor{codekw}{RGB}{140,40,140}
\definecolor{codecm}{RGB}{110,120,130}
\definecolor{codestr}{RGB}{30,120,60}

% "external link" glyph used in the JOSS margin block
\newcommand{\ExternalLink}{%
  \raisebox{-0.1ex}{\scriptsize$\nearrow$}}

% shorthand used throughout the text
\newcommand{\je}{\textsf{Jexpresso}}

% author/affiliation spacing
\renewcommand\Authfont{\normalsize}
\renewcommand\Affilfont{\small\itshape}
\setlength{\affilsep}{0.6em}

% tighter lists
\newcommand{\tightlist}{%
  \setlength{\itemsep}{0pt}\setlength{\parskip}{0pt}}

% Julia code listings
\lstdefinelanguage{Julia}{
  morekeywords={function,end,for,if,else,elseif,return,using,import,
                struct,mutable,const,module,in,where,let,do,begin},
  sensitive=true,
  morecomment=[l]{\#},
  morestring=[b]",
}
\lstset{
  language=Julia,
  basicstyle=\ttfamily\scriptsize,
  keywordstyle=\color{codekw}\bfseries,
  commentstyle=\color{codecm}\itshape,
  stringstyle=\color{codestr},
  backgroundcolor=\color{codebg},
  frame=single,
  framerule=0pt,
  xleftmargin=6pt,
  framexleftmargin=6pt,
  aboveskip=6pt,
  belowskip=4pt,
  columns=fullflexible,
  keepspaces=true,
  showstringspaces=false,
  breaklines=true,
}

% Bibliography environment produced by the JOSS/pandoc CSL pipeline.
% It must be a genuine list environment: the entries arrive as
% \bibitem[\citeproctext]{ref-key}, and \bibitem expands to \item here.
\newenvironment{CSLReferences}[2]{%
  \small
  \begin{list}{}{%
    \setlength{\leftmargin}{1.5em}%
    \setlength{\itemindent}{-1.5em}%
    \setlength{\labelwidth}{0pt}%
    \setlength{\labelsep}{0pt}%
    \setlength{\listparindent}{0pt}%
    \setlength{\itemsep}{0.18em}%
    \setlength{\parsep}{0pt}%
    \setlength{\topsep}{0.3em}%
  }%
}{\end{list}}
\newcommand{\citeproctext}{}
\renewcommand{\bibitem}[2][]{\item}


\title{Jexpresso.jl: An Extensible Julia Software for the
Solution of Partial Differential Equations on CPUs \& GPUs with Minimal User
Intervention}

        \author[1, 2]{Simone Marras}
        \author[3]{Yassine Tissaoui}
        \author[1]{Hang Wang}

      \affil[1]{Department of Mechanical \& Industrial Engineering, New
Jersey Institute of Technology, Newark, NJ, USA}
      \affil[2]{Department of Mathematical Sciences, New Jersey
Institute of Technology, Newark, NJ, USA}
      \affil[3]{Department of Mathematics, University of Wisconsin,
Madison, WI, USA}
  \date{\vspace{-5ex}}

\begin{document}
\maketitle

\marginpar{
  \sffamily\small

  {\bfseries DOI:} \href{https://doi.org/}{\color{linky}{}}

  \vspace{2mm}

  {\bfseries Software}
  \begin{itemize}
    \setlength\itemsep{0em}
    \item \href{https://github.com/openjournals/joss-reviews/issues/}{\color{linky}{Review}} \ExternalLink
    \item \href{https://github.com/smarras79/Jexpresso/}{\color{linky}{Repository}} \ExternalLink
    \item \href{https://github.com/smarras79/Jexpresso/}{\color{linky}{Archive}} \ExternalLink
  \end{itemize}

  \vspace{2mm}

  {\bfseries Submitted:} \\
  {\bfseries Published:}

  \vspace{2mm}
  {\bfseries License}\\
  Authors of papers retain copyright and release the work under a Creative Commons Attribution 4.0 International License (\href{https://creativecommons.org/licenses/by/4.0/}{\color{linky}{CC BY 4.0}}).
}

\section{Summary}\label{summary}

\je~is an open-source Julia package for solving systems of conservation laws in
one, two, and three spatial dimensions. Its guiding idea is that a user should be
able to add a new physical problem by writing down what they would write on
paper---a solution vector, its fluxes, its sources, and its boundary
conditions---and nothing else. Everything downstream, namely the spatial
discretization, stabilization, time integration, parallel decomposition, and the
choice of CPU or GPU execution, is supplied by the package and is deliberately
invisible to the problem definition.

The default spatial discretization is a continuous spectral element method (SEM)
on Legendre--Gauss--Lobatto nodes at arbitrary polynomial order, which is most
advantageous at fourth order and above, with a design that allows finite
difference, finite volume, or discontinuous Galerkin operators to be added
non-invasively. Time integration is delegated to the DifferentialEquations.jl
ecosystem (Rackauckas and Nie 2017), giving access to explicit and implicit
schemes without any solver code in \je~itself. Distributed-memory parallelism
uses MPI.jl (Byrne et al.\ 2021) through the Gridap.jl ecosystem (Badia and
Verdugo 2020), with meshes read via GridapGmsh.jl and adaptive refinement and
load balancing handled by p4est (Schlottke-Lakemper et al.\ 2023). Threading and
GPU offload are expressed once and dispatched to either target through
KernelAbstractions.jl and JACC.jl (Valero-Lara et al.\ 2024).

\section{Statement of Need}\label{statement-of-need}

Systems of conservation laws underpin computational fluid dynamics, atmospheric
and ocean modeling, radiative transfer, and astrophysics. Spectral element
methods deliver high accuracy per degree of freedom, but their implementation
complexity has historically kept them out of reach of researchers who are
domain experts rather than numerical analysts. Mature frameworks such as
Nek5000 (Fischer et al.\ 2008) and SPECFEM3D (Komatitsch and Tromp 2002) are
written in Fortran or C++, which raises the cost of prototyping a new equation
set. Julia packages for related problems---ClimateMachine.jl (Climate Modeling
Alliance 2019) and TrixiAtmo.jl (Ranocha et al.\ 2024)---target specific
equation sets or use discontinuous Galerkin rather than continuous spectral
elements.

\je~addresses the gap between ``a framework that solves my equations'' and ``a
framework I can teach my equations''. Adding a system requires no modification of
any solver: a self-contained problem directory supplies the initial condition,
fluxes, sources, boundary conditions, and the conserved-to-primitive mapping, and
that same directory then runs in 1D, 2D or 3D, in serial or under MPI, on CPU or
GPU, with any time integrator in the ecosystem. The repository ships roughly one hundred such directories across nine
equation families---from advection--diffusion and Burgers to compressible Euler,
elliptic and Helmholtz problems, ideal MHD, radiative transfer, shallow water
and acoustics---serving both as regression tests and as templates.

Two capabilities are, to our knowledge, not available together elsewhere in the
Julia ecosystem. First, Laguerre-based semi-infinite spectral elements
(Tissaoui et al.\ 2024) let a bounded computational domain be extended to
infinity, removing spurious reflections at artificial boundaries---essential for
atmospheric and oceanic problems. Second, a residual-based dynamic
sub-grid-scale model (Marras et al.\ 2015) sets the artificial viscosity from the
local residual of the governing equations rather than from a tuned constant, so
dissipation appears at shocks and under-resolved features and vanishes where the
solution is smooth. Because it reads only the residual, it applies unchanged to
whatever system the user defined; alongside it \je~offers Smagorinsky, Vreman and
WALE closures, artificial viscosity, and Boyd--Vandeven spectral filtering.

\section{Design: equations as fluxes and sources}\label{design}

\je~is built around the general conservation law
\begin{equation}
\label{eq:CL}
\left(\delta\,\partial_t {\bf q}\right) + \sum_{i=1}^{\rm nsd}\nabla\cdot{{\bf F}_i({\bf q})} = \mu\nabla^2{\bf \hat{q}} + {\bf S}({\bf q}) + ~{\rm b.c.},
\end{equation}
where $\nabla = [\partial_{x_i}]$, $\bf q$ is the solution vector, ${\bf F}_i$
and $\bf S$ are the flux and source vectors, and $\bf\hat q$ are the variables
to be diffused. The time derivative is parenthesized because the system need not
be time dependent, which is how elliptic and Helmholtz problems are expressed.

Any problem in the form of (\ref{eq:CL}) is defined by populating $\bf q$,
${\bf F}_i$, $\bf S$, $\bf\hat q$ and the boundary conditions. Two-dimensional
transport, $\partial_t\rho + \partial_x(\rho u) + \partial_y(\rho v) = 0$, is
the whole of the listing below.

\begin{lstlisting}[caption={Flux function for 2D scalar advection.}]
function user_flux!(F, G, q, args...)
    # q = [rho];  F = [u * rho];  G = [v * rho]
    u = 0.5   # x-velocity [m/s]
    v = 1.0   # y-velocity [m/s]
    F[1] = u * q[1]
    G[1] = v * q[1]
end
\end{lstlisting}

\noindent
The compressible Euler equations with gravity and tracers, the quasi-1D Euler
equations of a Laval nozzle, ideal magnetohydrodynamics with divergence
cleaning, and the five-dimensional radiative transfer equation are expressed the
same way, differing only in the length of $\bf q$ and the content of
${\bf F}_i$ and $\bf S$. This mirrors the separation of concerns familiar from
\texttt{OpenFOAM}, except that \je~asks for flux and source components rather
than the strong form.

\section{Key Features}\label{key-features}

\begin{itemize}
\tightlist
\item \textbf{Dimensional flexibility:} 1D, 2D and 3D from a single problem definition.
\item \textbf{High-order accuracy:} continuous spectral elements at arbitrary order; Laguerre semi-infinite elements for unbounded domains.
\item \textbf{Stabilization:} residual-based dynamic SGS, Smagorinsky, Vreman, WALE, artificial viscosity, spectral filters.
\item \textbf{Parallelism and adaptivity:} MPI decomposition with p4est adaptive refinement and load balancing.
\item \textbf{Portability:} one set of kernels for threaded CPU and GPU execution.
\item \textbf{Interoperability:} Gridap.jl meshing, Gmsh input, VTK and PNG output, and a CI suite of analytical and benchmark tests.
\end{itemize}

\je~has been applied by the authors to compressible and incompressible
Navier--Stokes flows, wave propagation, 3D radiative transfer, gravity waves,
shallow cumulus convection, and turbulent atmospheric boundary layers (Tissaoui
et al.\ 2025; Marras et al.\ 2025), and is intended for researchers developing
parameterizations, exploring numerical methods, or coupling problems across
disciplines.

\section{Acknowledgements}\label{acknowledgements}

We acknowledge the continuous support and input of Samuel Stechmann (University
of Wisconsin, Madison) on the extension of \je~beyond fluid dynamics, and
discussions with Hendrik Ranocha (Johannes Gutenberg University Mainz) and the
other authors of Trixi.jl.

\section*{References}\label{references}
\addcontentsline{toc}{section}{References}

\protect\phantomsection\label{refs}
\begin{CSLReferences}{1}{1}

\bibitem[\citeproctext]{ref-badia2020gridap}
Badia, Santiago, and Francesc Verdugo. 2020. {``Gridap: An Extensible
Finite Element Toolbox in {J}ulia.''} \emph{Journal of Open Source
Software} 5 (52): 2520. \url{https://doi.org/10.21105/joss.02520}.

\bibitem[\citeproctext]{ref-byrne2021mpi}
Byrne, Simon, Lucas C Wilcox, and Valentin Churavy. 2021. {``MPI. Jl:
Julia Bindings for the Message Passing Interface.''} \emph{Proceedings
of the JuliaCon Conferences} 1: 68.

\bibitem[\citeproctext]{ref-clima2019}
Climate Modeling Alliance. 2019. \emph{{ClimateMachine.jl}: A
Julia-Based Climate Model}.
\href{https://github.com/CliMA/ClimateMachine.jl}{Https://github.com/CliMA/ClimateMachine.jl}.

\bibitem[\citeproctext]{ref-fischer2008nek5000}
Fischer, Paul F., James W. Lottes, and Stefan G. Kerkemeier. 2008.
{``Nek5000: Open Source Spectral Element {CFD} Solver.''} \emph{Argonne
National Laboratory, Mathematics and Computer Science Division}.
\url{https://nek5000.mcs.anl.gov}.

\bibitem[\citeproctext]{ref-komatitsch2002spectral}
Komatitsch, Dimitri, and Jeroen Tromp. 2002. {``Spectral-Element
Simulations of Global Seismic Wave Propagation -- {I}. {V}alidation.''}
\emph{Geophysical Journal International} 149: 390--412.
\url{https://doi.org/10.1046/j.1365-246X.2002.01653.x}.

\bibitem[\citeproctext]{ref-marras2015dynsgs}
Marras, Simone, Murtazo Nazarov, and Francis X. Giraldo. 2015.
{``Stabilized High-Order {G}alerkin Methods Based on a Parameter-Free
Dynamic {SGS} Model for {LES}.''} \emph{Journal of Computational
Physics} 301: 77--101. \url{https://doi.org/10.1016/j.jcp.2015.07.034}.

\bibitem[\citeproctext]{ref-marras2025jexpresso}
Marras, Simone, Yassine Tissaoui, Hang Wang, and Samuel Stechmann. 2025.
\emph{{JEXPRESSO} V0.1: A {J}ulia-Language, User-Friendly, Multi-Physics
Parallel Solver for the Solution of Conservation Laws on {CPU}s and
{GPU}s}. UNAM.

\bibitem[\citeproctext]{ref-rackauckas2017differentialequations}
Rackauckas, Christopher, and Qing Nie. 2017.
{``{DifferentialEquations.jl} -- a Performant and Feature-Rich Ecosystem
for Solving Differential Equations in {J}ulia.''} \emph{Journal of Open
Research Software} 5 (1): 15. \url{https://doi.org/10.5334/jors.151}.

\bibitem[\citeproctext]{ref-ranocha2024trixiatmo}
Ranocha, Hendrik, Michael Schlottke-Lakemper, Andrew R. Winters, Erik
Faulhaber, Jesse Chan, and Gregor Gassner. 2024. \emph{{TrixiAtmo.jl}:
Atmospheric Modeling with {T}rixi.jl Using a Discontinuous {G}alerkin
Spectral Element Method}.
\href{https://github.com/trixi-framework/TrixiAtmo.jl}{Https://github.com/trixi-framework/TrixiAtmo.jl}.

\bibitem[\citeproctext]{ref-ranocha2023mpi}
Schlottke-Lakemper, Michael, Hendrik Ranocha, and Joshua Lampert. 2023.
\emph{{P4est.jl}: {J}ulia Bindings for the p4est Library}.
\href{https://github.com/trixi-framework/P4est.jl}{Https://github.com/trixi-framework/P4est.jl}.

\bibitem[\citeproctext]{ref-tissaoui2024efficient}
Tissaoui, Yassine, James F. Kelly, and Simone Marras. 2024. {``Efficient
Spectral Element Method for the {E}uler Equations on Unbounded
Domains.''} \emph{Applied Mathematics and Computation} 487: 129080.
\url{https://doi.org/10.1016/j.amc.2024.129080}.

\bibitem[\citeproctext]{ref-tissaoui2025parcfd}
Tissaoui, Yassine, Samuel Stechmann, Hang Wang, and Simone Marras. 2025.
\emph{Adaptive Mesh Refinement as a Pathway to Including Realistic
Radiation Models in Numerical Simulations of the Atmosphere}. UNAM.

\bibitem[\citeproctext]{ref-JACC}
Valero-Lara, Pedro, William F Godoy, Het Mankad, et al. 2024. {``{JACC:
Leveraging HPC Meta-Programming and Performance Portability with the
Just-in-Time and LLVM-based Julia Language}.''} \emph{Proceedings of the
SC '24 Workshops of the International Conference on High Performance
Computing, Network, Storage, and Analysis}.
\url{https://doi.org/10.1109/SCW63240.2024.00245}.

\end{CSLReferences}

\end{document}
